\begin{abstract}
Deploying electrocardiogram (ECG) classifiers across devices with different
memory and bandwidth budgets conventionally requires training and validating a
separate model for each target. Matryoshka Representation Learning (MRL)
offers an alternative: a single encoder whose embedding can be truncated to any
prefix length, so that one artefact serves every operating point. We present
the first systematic study of nested representation learning for biosignals,
and---in contrast to the promotional framing common to compression
papers---we test how far the resulting claims actually hold.
%
On PTB-XL ($n$=2{,}158 held-out recordings, five diagnostic superclasses), a
single Inception1D encoder trained with a six-level nesting objective attains a
macro AUC of \R{IncAUCsmall} at $d$=16 (64\,B per embedding) and
\R{IncAUClarge} at $d$=512 (2\,KB), across five independent training seeds. A
paired bootstrap test shows the two operating points are \R{EquivVerdict}
within a pre-registered margin of $\pm$\R{EquivMargin} AUC
($\Delta$=\R{EquivDiff}, 95\%~CI [\R{EquivLo}, \R{EquivHi}]), establishing
dimension invariance as a statistical result rather than an eyeballed one.
%
We add the three controls the compression literature usually omits.
\emph{First}, fixed-dimension baselines trained with the identical backbone,
schedule and seeds, which is the only comparison that isolates the effect of
nesting. \emph{Second}, measured rather than extrapolated computation: on real
hardware, end-to-end latency varies by \R{LatencySpread}\,ms across the entire
32-fold range of $d$, because the encoder executes in full regardless of the
truncation point. MRL therefore reduces embedding storage, transmission
volume and the number of maintained model artefacts---not inference compute,
a distinction the deployment literature routinely blurs. \emph{Third},
representation probing that tests whether the embedding really organises
coarse-to-fine: ridge probes recovering eleven physiological descriptors from
each prefix yield a rank correlation of \R{HierRho} ($p$=\R{HierP}) between
hypothesised and measured saturation order, so the intuitive
``rhythm-first, morphology-later'' account is \R{HierVerdict}.
%
We further evaluate zero-shot transfer to three external corpora, graceful
degradation under reduced-lead acquisition, and robustness to
wearable-characteristic artefacts at controlled SNR. We conclude that nested
embedding compression is a robust and useful property of ECG representations,
while deployment on wearable hardware remains an inference from clinical
resting-ECG evidence rather than a validated result.
\end{abstract}
